\documentclass[../DoAn.tex]{subfiles}
\begin{document}

\section{Đặt vấn đề}
\label{section:1.1}
Pháp luật điều chỉnh hầu hết các quan hệ trong đời sống, từ việc ký hợp đồng lao động, thuê nhà, mua bán tài sản cho đến xử phạt vi phạm giao thông hay chia thừa kế. Tuy vậy, việc tìm được đúng quy định áp dụng cho một tình huống cụ thể không đơn giản. Hệ thống văn bản quy phạm pháp luật của Việt Nam có số lượng rất lớn, được ban hành bởi nhiều cơ quan ở nhiều cấp, thường xuyên được sửa đổi, bổ sung hoặc thay thế, và được viết bằng ngôn ngữ chuyên môn khác xa cách người dân mô tả vấn đề của mình. Một người lao động muốn biết công ty có được giữ bản gốc căn cước của mình hay không sẽ không biết rằng câu trả lời nằm ở Điều 17 Bộ luật Lao động năm 2019, dưới tiêu đề ``Hành vi người sử dụng lao động không được làm khi giao kết, thực hiện hợp đồng lao động''.

Để giảm độ phức tạp này, Nhà nước đã thực hiện công tác pháp điển hoá theo Pháp lệnh Pháp điển hệ thống quy phạm pháp luật năm 2012 \cite{plpd2012}. Kết quả là Bộ Pháp điển điện tử \cite{phapdienportal}, trong đó các quy phạm pháp luật còn hiệu lực được sắp xếp theo 45 chủ đề, 306 đề mục, chương và điều, kèm liên kết giữa các điều có liên quan và chú thích văn bản gốc của từng điều. Đây là nguồn dữ liệu chính thống, có cấu trúc tốt, nhưng cách khai thác chủ yếu vẫn là duyệt cây chủ đề hoặc tìm kiếm theo từ khoá. Người dùng vì vậy vẫn phải tự biết mình cần tìm thuật ngữ gì, tự đọc hàng chục điều luật và tự đối chiếu với hoàn cảnh của mình.

Sự xuất hiện của các mô hình ngôn ngữ lớn (LLM) \cite{brown2020gpt3} mở ra khả năng hỏi đáp pháp luật bằng ngôn ngữ tự nhiên. Tuy nhiên, khi được dùng trực tiếp, LLM có hiện tượng ảo giác (hallucination) \cite{ji2023hallucination}: mô hình có thể đưa ra câu trả lời trôi chảy nhưng dẫn sai số điều, bịa số hiệu văn bản hoặc áp dụng quy định đã hết hiệu lực. Trong lĩnh vực pháp lý, một câu trả lời sai nhưng nghe có vẻ chắc chắn còn nguy hiểm hơn việc không có câu trả lời, vì người dùng có thể dựa vào đó để ký hợp đồng, khiếu nại hoặc từ bỏ quyền lợi của mình.

Nhu cầu không chỉ dừng ở việc hỏi đáp. Trước khi ký một hợp đồng mua bán, thuê nhà hay hợp đồng lao động, người dân thường không có điều kiện thuê luật sư rà soát, trong khi các điều khoản bất lợi hoặc trái luật như giữ giấy tờ tuỳ thân, phạt vi phạm vượt mức luật cho phép hay thoả thuận không đóng bảo hiểm xã hội vẫn xuất hiện phổ biến. Một hệ thống có thể đọc văn bản, chỉ ra điều khoản có nguy cơ trái luật và dẫn đúng điều luật làm căn cứ sẽ giúp người dân tự bảo vệ quyền lợi trước khi xảy ra tranh chấp.

Bài toán đặt ra vì vậy là xây dựng một hệ thống hỏi đáp và rà soát văn bản pháp luật tiếng Việt vừa dễ dùng như một cuộc trò chuyện, vừa đáng tin cậy như một tài liệu tra cứu: mỗi kết luận phải truy được về điều luật cụ thể, và hệ thống phải biết nói ``chưa đủ căn cứ'' thay vì suy diễn. Nếu giải quyết được, bài toán mang lại lợi ích cho người dân và doanh nghiệp nhỏ không có bộ phận pháp chế, cho cán bộ tư pháp cơ sở cần một công cụ tra cứu nhanh, và cách tiếp cận này có thể áp dụng cho các lĩnh vực có kho tri thức chính thống tương tự như y tế, thuế hay quy chế nội bộ của tổ chức.

\section{Mục tiêu và phạm vi đề tài}
\label{section:1.2}
Các sản phẩm hiện có để tra cứu pháp luật có thể chia thành ba nhóm. Nhóm thứ nhất là các cổng thông tin chính thống như Bộ Pháp điển điện tử và Cơ sở dữ liệu quốc gia về văn bản pháp luật; chúng cung cấp dữ liệu đầy đủ và tin cậy nhưng chỉ hỗ trợ duyệt và tìm kiếm theo từ khoá. Nhóm thứ hai là các trang thư viện pháp luật thương mại, bổ sung tình trạng hiệu lực, văn bản liên quan và dịch vụ tư vấn, nhưng phần lớn tính năng chuyên sâu yêu cầu trả phí và cách tra cứu vẫn dựa trên từ khoá. Nhóm thứ ba là các trợ lý AI đa năng; chúng hiểu câu hỏi bằng ngôn ngữ tự nhiên rất tốt nhưng không được gắn với một kho văn bản pháp luật Việt Nam cập nhật, không trích dẫn được điều luật có thể kiểm chứng và dễ nhầm lẫn văn bản cũ với văn bản mới. Phân tích chi tiết các nhóm sản phẩm này được trình bày trong Chương \ref{chapter:Related_works}.

Từ các hạn chế trên, đồ án đặt mục tiêu xây dựng một hệ thống hỏi đáp pháp luật Việt Nam trên dữ liệu Bộ Pháp điển, trong đó mọi câu trả lời và cảnh báo đều có căn cứ là điều luật thực sự tra được. Hệ thống hướng tới bốn nhóm chức năng chính: (i) hỏi đáp pháp luật bằng ngôn ngữ tự nhiên, hỗ trợ cả câu hỏi ngắn lẫn mô tả tình huống dài, có trích dẫn điều luật và cho phép xem toàn văn căn cứ; (ii) rà soát văn bản pháp lý do người dùng tải lên như hợp đồng mua bán, thuê nhà, lao động, trả về tóm tắt, ý chính và các cảnh báo kèm căn cứ pháp luật; (iii) cá nhân hoá qua bộ nhớ dài hạn có kiểm soát quyền riêng tư, quản lý hội thoại, hỏi bằng giọng nói và nghe câu trả lời; và (iv) công cụ quản trị cho phép cập nhật văn bản mới, theo dõi hiệu lực văn bản và đánh giá chất lượng hệ thống.

Phạm vi của đồ án giới hạn ở dữ liệu Bộ Pháp điển gồm 66.706 điều luật, cùng các văn bản được quản trị viên bổ sung qua công cụ nhập liệu, và các nguồn web pháp luật có độ tin cậy cao được dùng để kiểm tra quy định mới. Hệ thống đóng vai trò công cụ tham khảo, không thay thế ý kiến tư vấn của luật sư; mọi câu trả lời đều kèm lưu ý này. Đồ án không xây dựng mô hình ngôn ngữ riêng mà sử dụng các mô hình qua giao diện lập trình ứng dụng (API) của OpenAI, tập trung vào thiết kế luồng truy xuất, tác tử và cơ chế kiểm chứng căn cứ.

\section{Định hướng giải pháp}
\label{section:1.3}
Đồ án giải quyết bài toán theo hướng sinh văn bản tăng cường truy xuất (RAG) \cite{lewis2020rag}, mở rộng bằng đồ thị tri thức (Graph RAG) \cite{edge2024graphrag} và tác tử AI có khả năng gọi công cụ theo mô hình ReAct \cite{yao2023react}. RAG được chọn vì nó buộc mô hình đọc điều luật trước khi trả lời, nhờ đó câu trả lời truy được về nguồn và có thể cập nhật kiến thức mà không cần huấn luyện lại mô hình. Đồ thị tri thức được chọn vì Bộ Pháp điển vốn có cấu trúc phân cấp và quan hệ liên quan giữa các điều; khai thác các quan hệ này giúp tìm ra những điều luật mà tìm kiếm theo độ tương đồng ngữ nghĩa bỏ sót. Tác tử được chọn thay cho một bước truy xuất cố định vì câu hỏi pháp lý rất đa dạng: có câu chỉ cần một điều luật, có câu cần mở rộng sang điều liên quan, có câu về mức phạt cần kiểm tra văn bản mới nhất trên web.

Giải pháp cụ thể gồm năm thành phần. Thứ nhất, một tầng tiền xử lý phân loại yêu cầu để trả lời nhanh câu chào hỏi, từ chối nội dung rác và nhận diện người dùng có cảm xúc tiêu cực. Thứ hai, một tác tử gọi công cụ gồm tìm kiếm lai kết hợp vector ngữ nghĩa trong Milvus với BM25 tiếng Việt qua hợp nhất thứ hạng, tự bổ sung điều được dẫn chiếu và khoản ngoại lệ tra từ đồ thị điều luật, mở rộng theo đồ thị, đọc toàn văn điều luật và tìm kiếm web có lọc độ tin cậy nguồn; tác tử bị buộc kiểm tra web khi câu hỏi liên quan mức phạt, mức tiền hoặc khi kết quả trong Pháp Điển yếu. Thứ ba, tính năng rà soát văn bản với cơ chế kiểm chứng căn cứ: cảnh báo chỉ được đưa ra khi có điều luật thực sự tra được kèm đoạn trích nguyên văn khớp với nội dung điều luật. Thứ tư, bộ nhớ dài hạn theo hướng Mem0 \cite{chhikara2025mem0} có lọc dữ liệu nhạy cảm và cho phép người dùng xoá toàn bộ. Thứ năm, công cụ quản trị nhập văn bản mới qua khu chờ duyệt và quản lý hiệu lực văn bản.

Đóng góp chính của đồ án gồm phương pháp truy xuất lai cho văn bản pháp lý tiếng Việt, đồ thị dẫn chiếu với kiểm tra ngoại lệ, cơ chế kiểm chứng căn cứ cho tác vụ rà soát văn bản, bộ nhớ dài hạn bảo vệ quyền riêng tư và công cụ quản trị dữ liệu văn bản có theo dõi hiệu lực. Về kết quả, hệ thống đã được xây dựng hoàn chỉnh gồm giao diện web và máy chủ dịch vụ, vượt qua 372 kiểm thử đơn vị; trên bộ 50 câu hỏi đánh giá, truy xuất lai nâng P@10 từ 0,368 lên 0,544 và MRR từ 0,280 lên 0,471 so với cách truy xuất ban đầu.

\section{Bố cục đồ án}
\label{section:1.4}
Phần còn lại của báo cáo đồ án tốt nghiệp này được tổ chức như sau.

Chương \ref{chapter:Related_works} trình bày kết quả khảo sát các sản phẩm tra cứu pháp luật hiện có và nhu cầu của người dùng, từ đó xác định các chức năng cần xây dựng. Chương này mô tả biểu đồ use case tổng quát và các biểu đồ phân rã cho bốn nhóm chức năng hỏi đáp, rà soát văn bản, tài khoản và quản trị, hai quy trình nghiệp vụ chính, đặc tả chi tiết năm use case quan trọng và các yêu cầu phi chức năng về độ tin cậy của câu trả lời, hiệu năng, bảo mật và quyền riêng tư.

Chương \ref{chapter:Methodology} giới thiệu nền tảng lý thuyết và công nghệ sử dụng. Phần lý thuyết gồm mô hình ngôn ngữ lớn và cơ chế gọi công cụ, kiến trúc RAG, Graph RAG và tác tử ReAct, các kỹ thuật truy xuất như nhúng vector, chỉ mục HNSW, BM25 và hợp nhất thứ hạng RRF, cùng bộ nhớ dài hạn cho tác tử. Phần công nghệ phân tích lý do lựa chọn FastAPI, Next.js, Milvus, Neo4j, MySQL và MongoDB so với các lựa chọn thay thế.

Chương \ref{chapter:Experiment} trình bày phân tích thiết kế, triển khai và đánh giá hệ thống. Nội dung gồm kiến trúc phân tầng, thiết kế gói và lớp, thiết kế giao diện, thiết kế cơ sở dữ liệu trên bốn hệ lưu trữ, quy trình thu thập, làm sạch, lập chỉ mục vector, BM25 và xây dựng đồ thị tri thức từ dữ liệu Pháp Điển, các biểu đồ trình tự cho luồng hỏi đáp và rà soát văn bản, kết quả xây dựng ứng dụng, kiểm thử và mô hình triển khai.

Chương \ref{chapter:SolutionAndContribution} trình bày các giải pháp và đóng góp nổi bật của đồ án, gồm phương pháp truy xuất lai cho tiếng Việt pháp lý, tác tử tự kiểm tra quy định hiện hành, đồ thị pháp luật cho dẫn chiếu và kiểm tra ngoại lệ, cơ chế kiểm chứng căn cứ trong rà soát văn bản, bộ nhớ dài hạn bảo vệ quyền riêng tư và quy trình quản lý dữ liệu văn bản có theo dõi hiệu lực. Mỗi giải pháp được trình bày theo trình tự vấn đề, cách giải quyết và kết quả đo được.

Chương \ref{chapter:conclusion} tổng kết những gì đồ án đã làm được và chưa làm được, so sánh với các sản phẩm tương tự và đề xuất hướng phát triển, trong đó có mô hình lưu trữ điều luật theo thời gian để xử lý trọn vẹn các trường hợp sửa đổi, thay thế và hợp nhất văn bản.

\end{document}
